\section{Factorisation des lecteurs terminaux par naturalité, raffinement et passage à la limite}

Le point de départ matériel est l’unique constructeur simultané qui produit un seul continu enrichi doté de cinq opérations intrinsèques corrélatives. Les notations \(\mathfrak G_{\rm sp}\), \(\mathfrak G_{\rm sol}\), \(\mathfrak G_{\rm gau}\), \(\mathfrak G_{\rm coh}\) et \(\mathfrak G_{\rm det}\) désignent cinq vues complètes de ce même registre, reproduit avec ses fibres, relations, nombres, orientation, retenue, mémoire, frontière, parcours, multisection, survie, termes terminaux et lecteurs. Elles ne sont considérées ni comme cinq branches choisies ni comme cinq données indépendantes : on démontre leur compatibilité sur le même registre, leur naturalité sous raffinement, leur non-vacuité conjointe et leur passage à la limite, puis on les applique respectivement aux cinq problèmes. Aucune preuve de ce volume ne renvoie à une autre source externe pour compléter une définition, une identité ou une étape logique.

Chaque fermeture est exposée de manière autonome : définitions, démonstrations et compatibilités apparaissent avant sa conclusion classique. Dans chaque résolution, quatre couches distinctes sont mises en évidence :
\begin{enumerate}
  \item identité exacte à chaque niveau fini ;
  \item compatibilité des raffinements et préservation de l’histoire enrichie ;
  \item uniformité et passage à la limite, avec conservation du terme terminal ;
  \item reconnaissance de la conclusion classique dans le codomaine exprimé.
\end{enumerate}

La composition directrice des cinq fermetures conserve d’abord l’action corrélative unique, ne forme le continu qu’au passage à la limite, puis ouvre une vue pour chaque application :
\[
\begin{aligned}
 \mathcal E_k&\xrightarrow{\ \mathcal O_k\ }\mathcal Z_k
 \xrightarrow{\ \varprojlim_k\ }
 \bigl(\mathcal C_{\rm cont}^{\rm disc},\mathcal O_\infty\bigr),\\
 \bigl(\mathcal C_{\rm cont}^{\rm disc},\mathcal O_\infty\bigr)
 &\xrightarrow{\ \operatorname{View}_T\ }\mathfrak G_T^{\rm int}
 \xrightarrow{\ \operatorname{App}_{T,d_T}\ }\mathfrak G_T(d_T),\\
 \mathfrak G_T(d_T)&
 \xrightarrow{\widetilde{\mathcal A}_{T,d_T}}\mathfrak M_T
 \xrightarrow{\mathcal P_T}\mathcal C_T,
 \qquad T\in\{\mathrm{sp,sol,gau,coh,det}\}.
\end{aligned}
\]
Ici, $\mathcal O_k$ agit une seule fois sur l’état complet et conserve ensemble les cinq caractéristiques. $\operatorname{View}_T$ ne restreint pas le domaine ni ne sélectionne les histoires : il expose seulement les coordonnées d’une application au sein de la valeur complète de $\mathcal O_\infty$. La donnée classique $d_T$ n’entre que dans $\operatorname{App}_{T,d_T}$, après le continu, et ne participe pas à la génération de la caractéristique intrinsèque. La droite réelle est également une sortie ultérieure de $\operatorname{ev}_{\mathbb R}$, et non une prémisse de la chaîne.

Chaque chapitre distingue le domaine de départ, la construction, le certificat, la conclusion démontrée et ses réfutateurs exacts. Les développements de référence précédents construisent les cinq compositions dans leurs domaines respectifs ; les chapitres suivants déploient ces constructions et effectuent la reconnaissance finale. Aucune flèche n’est remplacée par un renvoi, une conclusion nominale ou une hypothèse classique.
